\documentclass[a4paper,11pt]{ctexart}
\usepackage{xcolor}
\usepackage{array}
\usepackage{booktabs}
\usepackage{hyperref}
\hypersetup{colorlinks=true,linkcolor=blue!70!black}
\linespread{1.35}

\newcommand{\draft}{\textcolor{orange!80!black}{\textbf{【待写】}}}

\title{\textcolor{blue!70!black}{\Huge\textbf{焊接技术模块 · 课程大纲}}}
\author{}
\date{}

\begin{document}
\maketitle
\thispagestyle{empty}

\section*{第零层 · 入门故事}

\begin{enumerate}
\item[\textbf{第1课}] \textbf{金属是怎么"粘"在一起的？} \draft \\
\textbf{内容：} 两块铁怎么变成一块铁？焊接的本质——加热到熔化→融为一体→冷却。对比胶水（物理黏附）和焊接（冶金结合）的本质区别。情景：如果摩天大楼的钢结构只用螺栓不用焊接会怎样？
\item[\textbf{第2课}] \textbf{为什么焊工要戴面罩？} \draft \\
\textbf{内容：} 电弧光中的紫外线和红外线——焊接面罩的滤光原理。和日食眼镜、电焊护目镜的对比。情景：一个焊工被弧光"闪了眼睛"是什么感觉？
\item[\textbf{第3课}] \textbf{太空里能焊接吗？} \draft \\
\textbf{内容：} 1969年苏联宇航员在太空做了第一次焊接实验——没有重力，融化的金属不会"流下来"。难点在哪里？太空焊接对未来空间站建设和卫星维修的意义。
\end{enumerate}

\section*{深入课程}

\begin{enumerate}
\item[\textbf{第1课}] \textbf{焊接冶金——金属熔化后发生了什么？} \draft \\
\textbf{内容：} 焊接热影响区（HAZ）的三个区域、奥氏体→马氏体相变、热裂纹与冷裂纹。情景：为什么焊接时不能着急冷却？
\item[\textbf{第2课}] \textbf{焊接工艺——从手工电弧焊到激光焊} \draft \\
\textbf{内容：} SMAW（手工电弧焊）、GMAW/MIG（气体保护焊）、TIG焊、SAW（埋弧焊）、激光焊和电子束焊。各自的热源、保护方式和适用范围。情景：特斯拉的一体化压铸vs传统焊接——两种制造哲学的对比。
\item[\textbf{第3课}] \textbf{焊接应力与变形} \draft \\
\textbf{内容：} 焊接残余应力的产生机制、变形控制方法（反变形、刚性固定、散热）。有限元预测焊接变形的原理。情景：为什么船体分段焊接时必须严格控制变形？
\item[\textbf{第4课}] \textbf{焊接检验——如何判断一个焊口合格？} \draft \\
\textbf{内容：} 无损检测（NDT）方法——超声波（UT）、X射线（RT）、磁粉（MT）、渗透（PT）。破坏性试验——拉伸、弯曲、冲击。情景：北京大兴机场的钢结构焊缝——每个都要100\%探伤吗？
\item[\textbf{第5课}] \textbf{特种焊接——水下、太空与核环境} \draft \\
\textbf{内容：} 水下湿法焊接与干法焊接、核电站反应堆压力容器的焊接、太空焊接的挑战。情景：福岛核电站的焊接修复——为什么是最困难的焊接工作之一？
\item[\textbf{第6课}] \textbf{焊接自动化——从机器人到智能焊接} \draft \\
\textbf{内容：} 焊接机器人的传感器系统（激光跟踪、视觉识别）、自适应焊接参数控制。情景：北京奔驰工厂的焊装车间——一个车身有5000多个焊点，全部由机器人完成，精度0.1mm。
\end{enumerate}

\end{document}
